% Figure for Calculus §97.
\begin{tikzpicture}[x=1.7cm, y=1.7cm, font=\scriptsize]
  \draw[-stealth, thin, black!70] (-1.85,0) -- (1.95,0) node[right] {$x$};
  \draw[-stealth, thin, black!70] (0,-1.35) -- (0,1.45) node[above] {$y$};
  \foreach \c in {0.5,1.5,2.5} \draw[figG!70!black, thin] (0,0) ellipse ({sqrt(\c)} and {sqrt(\c/2)});
  \draw[figG!70!black, very thick] (0,0) ellipse (1 and {sqrt(0.5)});
  \draw[figG!70!black, very thick] (0,0) ellipse ({sqrt(2)} and 1);
  \draw[figR, thick] (0,0) circle (1);
  \node[figG!70!black, anchor=south west, fill=white, inner sep=1pt] at ({sqrt(2)*cos(40)},{sin(40)}) {$c=2$};
  \node[figG!70!black, anchor=south, fill=white, inner sep=1pt] at (0.5,-0.6) {$c=1$};
  \node[figR, anchor=east] at (-1.25,1.2) {$x^2+y^2=1$}; \draw[figR, thin] (-1.25,1.2) -- ({cos(125)},{sin(125)});
  % parallel gradients at the touching points
  \foreach \p/\v in {{(1,0)}/{(1,0)}, {(-1,0)}/{(-1,0)}, {(0,1)}/{(0,1)}, {(0,-1)}/{(0,-1)}} {
    \fill \p circle (1.5pt);
  }
  \draw[-{Stealth[length=5pt]}, figR, thick] (1,0) -- (1.35,0);
  \draw[-{Stealth[length=5pt]}, figR, thick] (-1,0) -- (-1.35,0);
  \draw[-{Stealth[length=5pt]}, figR, thick] (0,1) -- (0,1.35);
  \draw[-{Stealth[length=5pt]}, figR, thick] (0,-1) -- (0,-1.3);
  \node[below right] at (1,0) {$(1,0)$};
  \node[above right] at (0,1) {$(0,1)$};
\end{tikzpicture}
